\section{创新点二\ 多感官互补启发的全域协同融合}
\subsection{问题和背景}
随着智慧城市基础设施持续向数字化、智能化和网络化方向发展，城市运行管理正逐步由局部、离散和被动的感知模式向全域、连续和智能的协同感知模式演进。在交通运行、公共安全、基础设施巡检、灾害应急以及无人系统自主作业等复杂城市场景中，感知对象具有空间分布广泛、动态变化迅速、环境干扰复杂以及信息来源多样等特点，对感知系统的覆盖范围、时空分辨率、实时响应能力和环境适应能力提出了更高要求。传统依赖单一传感器或单一感知节点的方式难以同时满足上述需求。视觉、雷达、红外、声学以及惯性等不同感知手段在观测机理和信息表达能力上存在天然差异，分别擅长捕获环境中的不同信息。例如，视觉传感器能够提供丰富的纹理、结构和语义信息，但容易受到光照变化、遮挡和运动模糊影响；雷达具有较强的全天候和远距离感知能力，但在空间分辨率和语义表达方面存在局限；惯性传感器能够提供高频、低延迟的运动信息，却难以长期避免累积误差。因此，复杂环境下的高质量感知并非依赖单一传感器能力的简单增强，而需要充分挖掘不同感知模态之间的能力互补关系，通过异构信息的协同表征与融合突破单一模态的感知边界。

进一步地，当城市感知系统从单一智能体向由大量终端、边缘节点和中心节点组成的分布式感知网络演进后，感知能力的提升又带来了新的系统性矛盾。一方面，不同空间位置的无人机、车辆、机器人、固定摄像头以及其他感知节点能够从不同视角和不同时间尺度观测同一环境，从而弥补单节点在视场范围、空间遮挡和局部感知能力上的不足；另一方面，节点数量和感知数据规模的快速增长，使原始图像、点云以及高维中间特征的大规模交换面临越来越严重的通信和计算压力。在有限带宽和有限计算资源条件下，并非所有感知信息都具有同等的传输价值：部分信息与当前任务高度相关，需要及时传递至关键节点，而大量冗余或低价值信息则无需持续传播。因此，如何从海量感知信息中识别与任务相关的高价值内容，并在跨时间、跨空间和跨节点的条件下进行有选择、高效率的信息传导，成为大规模协同感知进一步提升实时性和系统效率所必须解决的问题。

然而，即使通过多模态互补和跨节点协同获得了更加丰富的感知信息，多源观测本身仍不可避免地存在不确定性。复杂城市环境中的光照变化、遮挡、天气干扰、设备噪声、传感器漂移、通信异常以及不同传感器之间的认知偏差，都可能导致不同来源的信息存在质量差异、局部失真甚至相互冲突。在这种情况下，简单地增加感知来源并不能必然提升最终感知结果，反而可能将低质量甚至错误信息进一步传播和放大。传统融合方法通常依赖固定权重、统计一致性或特征级关联进行信息融合，对于“当前信息是否可信”“不同来源的信息应该采信多少”以及“当环境变化导致信息质量发生改变时如何动态调整”等问题缺乏充分建模。尤其当不同模态具有不同观测机理时，仅依据数据之间的一致性进行可靠性判断可能难以得到稳定结果。因此，在多源信息融合过程中，还需要进一步引入环境规律、任务知识、空间拓扑以及历史经验等先验信息，为信息可信性判断提供独立于当前观测的一致性依据，并据此实现融合过程的动态调节和持续校正。

上述问题与生物多感官系统的信息处理机制具有高度的对应关系。生物并不是通过单一感官获取完整环境信息，而是利用视觉、听觉、触觉等不同感官之间的优势互补形成更加全面的环境认知；面对持续输入的大量感知刺激，神经系统也不会对所有信息进行等价处理，而是根据当前任务和环境状态选择性地增强、抑制和传导重要信息，从而在有限神经资源下实现高效的信息处理；同时，生物感知还会结合已有经验、环境规律和先验知识，对当前观测进行解释和修正，以提高认知结果在复杂、不确定环境中的稳定性和可靠性。由此可以看到，感官互补解决的是“信息从哪里来、如何互补”的问题，选择传导解决的是“哪些信息值得传、如何高效传”的问题，而先验调节解决的是“哪些信息可信、如何动态校正”的问题。三者分别对应多源感知系统中的信息获取、信息流动和信息认知三个关键环节，为构建面向复杂城市环境的全域协同融合体系提供了重要启发。

基于此，本研究面向复杂城市场景下单一感知能力受限、分布式感知通信资源受限以及多源信息可信性不足等关键问题，借鉴生物多感官系统“优势互补—选择传导—先验调节”的信息处理机制，构建“多模态互补—跨时空传导—先验调节增强”的全域协同融合技术体系。首先，从感知源头出发，研究多模态感官之间的能力互补与异构信息融合机制，使不同传感器能够在统一感知框架下充分发挥各自优势，提升复杂环境中的感知完整性与精度；其次，从信息流动过程出发，研究面向任务价值的跨时空选择传导机制，在有限通信资源条件下识别并传递高价值感知信息，实现分布式节点之间的高效协同；最后，从融合结果出发，研究先验知识驱动的多源可信调节机制，根据环境规律和任务知识动态评估不同信息源的可靠性，对融合过程进行自适应调节和持续校正，从而提高全域感知结果的可靠性与鲁棒性。三个研究层次由底层感知信息获取逐步延伸至网络信息传导，再进一步作用于高层可信认知，形成从“感知互补”到“信息协同”再到“可信融合”的完整技术链路。

\textbf{多模态感官互补的异构融合。}在复杂城市环境中，无人机、移动机器人等智能终端需要依赖多种传感器实现连续、精准的自主定位与环境感知。然而，不同感知模态具有不同的物理观测机制和信息表达能力，单一传感器往往难以同时满足高动态、高精度和全天候感知需求。以自主定位为例，惯性测量单元（IMU）能够提供高频率、低延迟的运动状态信息，具有良好的短时动态响应能力，但其本质上属于相对运动感知，误差会随时间不断累积，难以长期提供稳定的绝对位置约束；RGB/D相机能够从环境中提取丰富的纹理、结构和语义信息，并结合视觉惯性里程计或视觉定位服务实现较高精度的定位，但在高速运动、剧烈姿态变化和复杂光照条件下，传统帧式相机容易产生运动模糊，其有限帧率也限制了对高速动态平台的高频状态更新能力；激光雷达和毫米波雷达具有较强的环境适应能力，但分别受到采样频率、空间分辨率、语义表达能力以及设备成本等因素的制约。因此，不同传感器并不存在绝对意义上的优劣，而是在时间分辨率、空间结构、语义信息和环境适应性等维度形成不同的能力优势。

现有无人机视觉定位系统通常采用“RGB相机+IMU+先验三维地图”的组合方式，由RGB相机提取环境视觉特征，IMU提供高频相对运动约束，三维地图提供绝对空间参考。这种组合能够在一般场景下获得较好的定位性能，但在高速无人机降落等高动态场景中仍面临明显瓶颈：高速运动容易造成RGB图像运动模糊，使传统视觉特征提取和匹配受到影响；视觉惯性定位通常只能维持在几十Hz量级，而二维图像与三维地图之间的特征匹配又带来较高计算开销，基于地图的绝对定位更新频率甚至可能低于1 Hz，难以满足飞控系统对高频状态反馈的需求。

事件相机为解决这一问题提供了新的感知能力。作为一种具有仿生特征的视觉传感器，事件相机以异步方式记录像素亮度变化，具有微秒级时间分辨率、低延迟和近乎无运动模糊等特点，尤其适合高速动态场景下的视觉感知。由此，事件相机、IMU和先验三维地图形成了具有明显互补性的感知组合：事件相机能够捕获高速动态视觉变化，IMU能够提供高频运动状态约束，而三维地图能够提供稳定的全局空间参考。EV-Pose正是在这一背景下，探索三类异构信息之间的协同融合机制，以充分发挥不同感知模态在时间、空间和全局参考维度上的互补优势。

然而，多模态传感器的简单叠加并不能自动形成有效的融合能力。事件流具有异步、稀疏和高时间分辨率等特征，与传统图像以及三维点云地图在数据结构和信息表达方式上存在显著差异，形成明显的2D–3D异构模态鸿沟；同时，高速运动会产生大量事件数据，如果不加选择地将全部事件输入后续计算过程，也会显著增加系统计算负担。因此，关键不在于简单增加传感器数量，而在于理解不同感知模态的能力边界与物理关联，并围绕任务需求建立深层次的协同表征与融合机制。这一问题构成了多模态感官互补异构融合研究的核心。


\subsection{研究内容与目标}

\subsection{研究方法}

\subsection{实验结果与分析}

\subsection{讨论与总结}

% 创新点二：多感官互补启发的全域协同融合。
% 针对智慧城市基础设施建设中感知手段单一、通信带宽受限、采信判据缺失等问题，借鉴生物多感官系统的感官互补、选择传导和先验调节机制，提出多模态感官互补、跨时空选择传导和先验知识调节增强的多源融合方法。通过异构模态优势互补提升观测与定位精度，通过高价值信息选择传导降低通信开销，通过先验知识调节实现多源观测的可信采信与动态校正，形成多模态互补-跨时空传导-先验调节增强的多源融合技术体系，显著提升复杂城市场景下全域感知的完备性、实时性和可靠性。

% 2.1多模态感官互补的异构融合方法
% 针对复杂城市场景中单一感知手段信息维度有限、异构模态观测特性差异显著，难以形成准确完备观测的问题，提出多模态感官互补的异构融合方法。融合事件视觉、惯性运动与三维空间结构等异构感知信息，构建跨模态时空特征匹配与运动感知层级融合机制，协同利用高时间分辨视觉、连续运动约束和三维空间结构信息，实现不同模态在时间、运动和空间维度的优势互补，增强复杂动态环境下异构观测的协同表征与融合定位能力。相关方法在真实场景定位实验中优于香港大学、美国特拉华大学、西班牙萨拉戈萨大学等团队的代表性方法，多模态融合定位误差降低30%以上，显著提升多源感知的准确性与观测完备性。
% 直接支撑论文：EV-Pose

% 2.2跨时空选择传导的多节点协同融合方法
% 针对城市大规模感知中多节点观测冗余与有限通信资源之间的矛盾，提出跨时空选择传导的多节点协同融合方法。建立感知质量驱动的节点调度与信息价值评价机制，利用节点间视角互补、空间位置关联和任务时序一致性，动态筛选高价值观测；结合事件触发、向量量化与熵编码等技术，压缩语义与统计冗余，降低低价值信息的重复传输；构建跨节点信息扩散与协同融合机制，实现关键信息在有限通信资源下的高效传导。相关方法在感知-通信权衡方面优于上海交通大学、上海人工智能实验室、北京航空航天大学等团队的代表性方法，在同等感知精度条件下通信带宽降低55%以上，显著提升大规模协同感知的实时性和通信效率。
% 直接支撑论文：Zeng Y, Li S, Li Z, et al. Communication-Efficient Collaborative Perception with Semantic and Statistical Compression[C]//Pattern Recognition and Computer Vision: 8th Chinese Conference, PRCV 2025, Proceedings, Part XI. 2025: 270-283.

% 2.3先验知识调节增强的多源可信融合方法
% 针对复杂城市场景中多源观测质量动态波动、可靠信息难判别、异常观测难校正的问题，提出先验知识调节增强的多源可信融合方法。融合感知可靠性以及语义、结构和运动等先验知识，构建多源观测的可信判别、选择采信与动态校正机制。基于可靠性评估与先验约束，动态判断不同来源观测的可信程度，对低可信、偏移和失真信息进行自适应筛选与校正，降低环境变化、域偏移和复杂运动对融合结果的干扰，实现由等权融合向可信采信、先验调节转变。相关方法在综合感知质量与空间重建性能方面优于卡内基梅隆大学、中南大学、北京航空航天大学等团队的代表性方法，空间感知误差降低39%，显著提升复杂动态环境下全域感知的准确性与可靠性。
% 直接支撑论文：Zhou N, Li Z, Man F, et al. QUIDS: Quality-Informed Incentive-Driven Multiagent Dispatching System for Mobile Crowdsensing[J]. IEEE Internet of Things Journal, 2026, 13(8): 15235-15249.
